\section{Background and Related Work}\label{sec:background}

\subsection{The SQLite WAL Contract}
SQLite's write-ahead log~\cite{sqlite_wal} provides snapshot reads with
a single-writer constraint. A new connection acquires the appropriate
shared or exclusive lock; readers see the most recent committed
checkpoint; the writer appends frames to the WAL file and commits by
writing a frame whose header marks end-of-transaction. The contract is
robust and easy to reason about, but it serializes every writer through
one lock. On a 64-thread workload mixing reads and writes, throughput
plateaus at the rate at which a single writer can fsync.

\subsection{ARIES, B-link Trees, and the Index Concurrency Problem}
ARIES~\cite{mohan1992aries} established the canonical pattern of
write-ahead logging with redo/undo records, a checkpoint that does not
require quiescing the database, and physiological logging that combines
page-level redo with logical undo. RedlineDB's WAL coordinator follows
this template.

For the index layer, the obvious analogue of ``serialize one writer''
is to lock the entire B-tree at the root. Lehman and Yao showed in
1981~\cite{lehman1981blink} that B-link trees can support concurrent
inserts, splits, and lookups by augmenting each internal node with a
right-link pointer and following the link when a search overshoots a
recently split node. Srinivasan and Carey~\cite{srinivasan1991bplus}
empirically validated that this discipline outperforms lock-coupling
on multi-core hardware. Graefe~\cite{graefe2010btree} surveys the
broader space of B-tree locking strategies. Bayer and
McCreight~\cite{bayer1972btree} are the original B-tree reference. We
adopt the B-link mental model: index leaves carry transactional
mark/unmark operations rather than physical deletions during a live
transaction, which lets the executor undo a partially-applied DML
without quiescing the index.

\subsection{MVCC and Snapshot Isolation}
Multi-version concurrency control~\cite{bernstein1983mvcc} decouples
readers from writers by keeping multiple versions of a row, indexed by
commit timestamps. Reed's 1978 thesis~\cite{reed1978naming} introduced
the per-version naming scheme that all modern MVCC engines descend from.
Berenson et al.~\cite{berenson1995ansi} cataloged the anomalies---write
skew, read-only-anomaly, lost update---that show up between read
committed and serializable, and Cahill et
al.~\cite{cahill2008ssi} introduced serializable snapshot isolation
(SSI) as a runtime discipline that detects dangerous structures. Kung
and Robinson~\cite{kung1981occ} are the canonical reference for
optimistic concurrency control.

RedlineDB currently implements snapshot isolation: a transaction sees
all rows committed at or before its read-CSN, and writes acquire row
locks rather than running a certification phase. SSI is future work
(Section~\ref{sec:discussion}). In practice, SI has been the production
isolation level of PostgreSQL~\cite{postgresql_mvcc,stonebraker1987postgres},
InnoDB~\cite{innodb_mvcc}, and Hekaton~\cite{diaconu2013hekaton} for
years; we follow the same convention. Gray's classic transaction-concept
paper~\cite{gray1981transaction} and the granularity-of-locks
report~\cite{gray1976granularity} frame why row-granularity locks
matter for the workloads where MVCC alone is not enough.

\subsection{Embedded Databases and Where RedlineDB Sits}
LMDB~\cite{chu2011lmdb} is the canonical memory-mapped embedded store:
single-writer, copy-on-write B$^+$-tree, no SQL. LevelDB~\cite{leveldb}
and RocksDB~\cite{rocksdb} are LSM-based key-value stores; they are
write-throughput optimized but expose no relational surface. In the
Rust ecosystem, sled~\cite{sled} (now in maintenance mode) and
redb~\cite{redb} are pure-Rust key-value stores; neither speaks SQL or
the SQLite ABI. libSQL~\cite{libsql} is Turso's open-contribution fork
of SQLite written in C and Rust; it preserves the SQLite WAL model and
its single-writer constraint by construction. DuckDB~\cite{raasveldt2019duckdb}
is the embedded analytical comparison point, but its design center is
column-store OLAP, not OLTP. HyPer~\cite{kemper2011hyper} and the fast
serializable MVCC follow-up~\cite{neumann2015fast} are closer in
spirit---OLTP, MVCC, in-memory---but live in a server context. RedlineDB
occupies the cell that none of these fill: embedded, SQL, MVCC,
multi-writer, SQLite-ABI-compatible.

\subsection{Determinism, Failpoints, and Crash Testing}
Failpoint discipline---compile-neutral injection points wired to a
text-driven schedule---was popularized by TiKV's
\texttt{fail-rs}~\cite{failrs}. SQLite's own testing
program~\cite{sqlite_testing} and the SQLLogicTest
harness~\cite{sqllogictest} set the durability bar; Jepsen reports
contextualize what storage engines can and cannot
promise~\cite{jepsen_sqlite}. FoundationDB~\cite{zhou2021foundationdb}
and TigerBeetle~\cite{tigerbeetle_sim} have shown how deterministic
simulation lets a small team certify a transactional kernel against
adversarial schedules; whitebox fuzzing~\cite{godefroid2008whitebox} and
property testing~\cite{proptest,claessen2000quickcheck} round out the
modern toolkit. RedlineDB's certification protocol
(Section~\ref{sec:method}) draws from these lineages. The Rust safety
story~\cite{klabnik2023rust,jung2021rustbelt} and the formalization of
``safe outside \texttt{unsafe}''~\cite{rust_unsafe_guidelines} are what
make a deterministic kernel feasible without C-style memory hazards.
Bw-trees~\cite{levandoski2013bwtree} and consensus
protocols~\cite{ongaro2014raft} are noted here as adjacent design
points we did not adopt.
